%---------------------------------------------%
%--------          MY RESUME          --------%
%---------------------------------------------%
\documentclass[letter,9pt]{article}

%% PACKAGES %%
\usepackage[T1]{fontenc} % better PDF fonts/encoding
\usepackage[english]{babel} % language support
\usepackage[utf8]{inputenc} % encoding
\usepackage{lmodern} % Latin Modern fonts
\usepackage[bottom=0.5in, top=0.5in, left=0.5in, right=0.5in]{geometry} % set all margins to 0.5 inches
\usepackage{multicol} % multiple columns
\usepackage{enumitem} % lists 
\usepackage{microtype} % better spacing

%% GLOBAL FORMATTING %%
\setlength{\parindent}{0pt} % Remove paragraph indentation globally

%% FORMATING %%
\pagestyle{empty} % No header, footers, or page number
\linespread{0.92}
\setlist{leftmargin=0.12in,itemsep=0pt,topsep=1pt,parsep=0pt} % tighter lists

\begin{document}
    %% HEADER %%
    \begin{center}
        {\Large \textbf{PONTHEA ZAHRAII}} \\
        {\small Software Engineer II $\bullet$ Robotics Perception, 3D Vision, ML Systems $\bullet$ US Citizen} \\
        {\small (949) 397-7414 $\bullet$ pontheazahraii@gmail.com $\bullet$ www.linkedin.com/in/pontheazahraii $\bullet$ www.github.com/pontheazahraii}
    \end{center}
    
    \vspace{3pt}
    %% EDUCATION %%
    \noindent \textbf{EDUCATION} \hrulefill
    \vspace{2.5pt}
    
    \noindent
    \begin{tabular*}{\textwidth}{@{\extracolsep{\fill}} l r}
        \textbf{Chapman University, Fowler School of Engineering} & Orange, CA \\
        Master of Science, Electrical Engineering and Computer Science & GPA 4.00 | May 2025 \\
        \textit{Thesis: Implementation of Residual Tandem Neural Networks in Inverse Photonic Design} & ~ \\
        \textit{Valedictorian, Graduation Keynote Speaker} & ~ \\[5pt]
        Bachelor of Science, Computer Science, Neuroscience Minor & GPA 3.89 | May 2024 \\
        \textit{Graduation Keynote Speaker} & ~ \\
    \end{tabular*}

    \vspace{5pt}
    %% TECHNICAL SKILLS %%
    \noindent \textbf{TECHNICAL SKILLS} \hrulefill
    \vspace{2.5pt}
    
    {\footnotesize
    \noindent
    \textbf{Languages:} Python, C++, SQL, Bash

    \textbf{Robotics \& 3D Perception:} 
    ROS 2, TF2, RGB-D Perception, LiDAR Processing,
    Point Cloud Processing, ICP Registration,
    Multi-View Fusion, Geometric Registration,
    Coordinate Frame Systems, Spatial Alignment,
    Volumetric Mapping (TSDF/NVBlox),
    Occupancy Reasoning, Scene Reconstruction,
    Surface Meshing, Target Point Generation,
    Sensor Fusion, Real-Time Perception Systems,
    RViz

    \textbf{Computer Vision \& ML:} 
    TensorFlow, OpenCV, Open3D, PCL,
    Object Detection, Segmentation, Classification,
    Feature Extraction, Depth Estimation,
    Geometric Computer Vision, Model Evaluation

    \textbf{Systems \& Infrastructure:} 
    Docker, Linux, Git, CI/CD (GitHub Actions),
    Google Cloud Platform (GCP),
    GPU Acceleration, Parallel Processing,
    MLOps, Real-Time Data Pipelines, \LaTeX
    }

    \vspace{5pt}
    %% PUBLICATIONS %%
    \noindent \textbf{PUBLICATIONS} \hrulefill
    \vspace{2.5pt}
    
    {\footnotesize
    \begin{itemize}[leftmargin=*,noitemsep,topsep=0pt]
        \item Estakhri, N.M., Zahraii, P., Kashanchi, S., Estakhri, N.M. (2024). ``Predictive residual neural networks for optical trapping of small particles.'' \textit{Optical Trapping and Optical Micromanipulation XXI}, SPIE, Vol. 13112, pp. 111--116.
        \item Zahraii, P., Kashanchi, S., Estakhri, N.M., Estakhri, N.M. (2025). ``Machine learning-driven prediction and inverse design of optical forces near gradient metasurfaces.'' \textit{Optical Trapping and Optical Micromanipulation XXII}, SPIE, PC135800J.
        \item Zahraii, P., Kashanchi, S., Estakhri, N.M., Estakhri, N.M. (2025). ``A Deep Learning Framework for Prediction and Inverse Design of Nanoscale Optical Forces near Gradient Metasurfaces.'' \textit{2025 Nineteenth International Congress on Artificial Materials for Novel Wave Phenomena (Metamaterials)}, IEEE, X--389.
        \item Zahraii, P.A. (2025). ``Implementation of Residual Tandem Neural Networks for Photonic Inverse Design.'' Master's Thesis, Chapman University.
        \item Zahraii, P., Kashanchi, S., Estakhri, N.M., Estakhri, N.M. (2025). ``Machine learning-assisted characterization of optical forces near gradient metasurfaces.'' \textit{SMT 2025}, APS.
    \end{itemize}
    }

    \vspace{5pt}
    %% PROFESSIONAL EXPERIENCE %%
    \noindent \textbf{PROFESSIONAL EXPERIENCE \& RESEARCH} \hrulefill
    \vspace{2.5pt}

    \noindent
    \begin{tabular*}{\textwidth}{@{\extracolsep{\fill}} l r}
        \textbf{Kazvu Labs}, Irvine, CA & November 2025 -- Present \\
    \end{tabular*}
    \noindent Software Engineer II \\
    \textit{Summary: Build robotics perception systems for real-time 3D scene understanding, mapping, reconstruction, and safety-oriented spatial reasoning.}
    {\footnotesize
    \begin{itemize}[noitemsep, topsep=0pt]
        \item Designed and developed end-to-end robotics perception pipelines using ROS 2, RGB-D sensors, LiDAR, Open3D, and PCL for real-time and offline 3D scene understanding.
        \item Built systems for point-cloud generation, segmentation, classification, target point generation, geometric reconstruction, and spatial reasoning from multi-sensor perception data.
        \item Implemented ICP-based registration, TF2 coordinate transforms, and multi-view geometric fusion workflows for scene alignment, spatial registration, and reconstruction.
        \item Developed volumetric mapping and safety-oriented perception systems using voxelized occupancy reasoning, TSDF/NVBlox concepts, intrusion detection, and online scene understanding.
        \item Built reusable C++ and Python perception libraries for filtering, visualization, feature extraction, registration, reconstruction, and large-scale point-cloud processing workflows.
        \item Integrated perception systems into ROS 2 robotics stacks with sensor ingestion pipelines, topic remapping, transform-tree management, and real-time streaming workflows.
        \item Optimized large-scale 3D processing workflows through GPU acceleration, parallelized processing, and reusable pipeline architecture.
        \item Maintained CI/CD workflows, Dockerized development environments, and engineering documentation while collaborating across robotics, software, and research teams.
    \end{itemize}
    }
    \vspace{5pt}

    \noindent
    \begin{tabular*}{\textwidth}{@{\extracolsep{\fill}} l r}
        \textbf{Hansji Corporation}, Irvine, CA & June 2024 -- November 2025 \\
    \end{tabular*}
    \noindent Machine Learning Engineer II \\
    \textit{Summary: Progressed from early RMS prototyping work to ownership of production ML, forecasting, and MLOps systems for revenue optimization.}
    {\footnotesize
    \begin{itemize}[noitemsep, topsep=0pt]
        \item Developed predictive models and decision algorithms for booking pace forecasting, dynamic pricing, pricing elasticity, and revenue optimization across a multi-property portfolio.
        \item Built the RMS from prototype algorithms into a production ML platform through implementation of ingestion, feature engineering, training, validation, deployment, and monitoring workflows.
        \item Designed scalable MLOps infrastructure and GCP integrations supporting reproducible experimentation, automated retraining, near real-time rate updates, and production pricing automation.
        \item Evaluated time-series, regression ensemble, and rule-based optimization approaches to improve forecasting accuracy and pricing decision quality.
        \item Partnered with BI analysts and leadership on feature analysis, benchmarking, A/B testing, and translation of model outputs into operational business decisions.
        \item Developed CI/CD workflows, engineering documentation, and reproducible ML tooling to improve maintainability and deployment consistency.
    \end{itemize}
    }
    \vspace{5pt}

    \noindent
    \begin{tabular*}{\textwidth}{@{\extracolsep{\fill}} l r}
        \textbf{Chapman University, Estakhri Lab}, Orange, CA & August 2023 -- August 2025 \\
    \end{tabular*}
    \noindent Researcher \\
    \textit{Summary: Conducted research on neural networks and inverse design for photonics, publishing peer-reviewed work.}
    {\footnotesize
    \begin{itemize}[noitemsep, topsep=0pt]
        \item Researched inverse design problems in photonics using residual tandem neural networks (forward + inverse models)
        \item Implemented residual blocks to stabilize training and improve accuracy for ill-posed inverse mappings
        \item Built data pipelines and training/validation workflows; tuned hyperparameters and evaluated generalization under noise
        \item Published peer-reviewed papers; contributed methods, experiments, and writing
    \end{itemize}
    }
\end{document}
